\uuid{zzBF}
\titre{Partie entière}
\auteur{bodin ; hanani ; mzari}
\organisation{exo7}
\licence{CC BY-NC-SA 4.0}
\datecreate{2020-01-20}
\niveau{L1}
\module{Analyse}
\chapitre{Propriétés de R}
\sousChapitre{Propriétés des nombres réels}
\difficulte{2}
\theme{Propriétés de nombres réels, partie entière}
\exo7qcmid{lille-1.reels.012}

\contenu{
\begin{question}
\qclassification{120.03}
\qtags{motcle=partie entière}
Notation : $E(x)$ désigne la partie entière du réel $x$.
Quelles sont les assertions qui caractérisent la partie entière ?
\begin{answers}
    \bad{$E(x)$ est le plus petit entier supérieur ou égal à $x$.}

    \good{$E(x)$ est le plus grand entier inférieur ou égal à $x$.}

    \bad{$E(x)$ est l'entier tel que $x \le E(x) < x+1$}

    \good{$E(x)$ est l'entier tel que $E(x) \le x < E(x)+1$}
\end{answers}
\begin{explanations}
$E(x)$ est le plus grand entier inférieur ou égal à $x$, ce qui se caractérise aussi par $E(x) \le x < E(x)+1$.
\end{explanations}
\end{question}
}
